% ORIGINAL RESEARCH DRAFT -- VERSION 2 (EDITORIAL REVISION), 23 NOVEMBER 2026
% Evidence base: PKgamerzop/pq-watermark-2.0, phase6 commit fc526ae62...
% and scientific corrections on paper-audit-fixes / paper-ieee-final.
% Version 2 changes wording, organization and framing only. All numerical
% results, denominators, equations, labels and figure placeholders from
% Version 1 are unchanged. See REVISION_NOTES.md for the author-approval list.
% WARNING: This is NOT a conference-specific camera-ready template.
% Authors must verify all content, authorship, AI disclosures and venue policies.
% To add original figures: place PDFs at the paths named in each \figslot call.
\documentclass[conference]{IEEEtran}
\IEEEoverridecommandlockouts
\usepackage{cite}
\usepackage{amsmath,amssymb}
\usepackage{graphicx}
\usepackage{booktabs}
\usepackage{array}
\usepackage{url}
\usepackage{xcolor}
\usepackage{balance}
\usepackage{microtype}

% A missing image is deliberately rendered as a labeled placeholder. Replace
% each one with an ORIGINAL diagram/plot generated from the project data.
\newcommand{\figslot}[3]{%
  \IfFileExists{#1}{\includegraphics[width=\linewidth]{#1}}{%
  \fbox{\begin{minipage}[c][#2][c]{0.94\linewidth}
     \centering\footnotesize\sffamily\bfseries FIGURE PLACEHOLDER\\[3pt]
     \normalfont\sffamily #3\\[3pt]
     \scriptsize Replace with an original, data-based figure.
  \end{minipage}}}}

\title{Post-Quantum-Signed Model Provenance Through Neural Weight Watermarking: An Empirical Study of Accuracy and Overwrite Resistance}

% AUTHOR CHECK: Names and affiliations below are taken from the project's
% existing conference source. Confirm authorship, order, spelling and whether
% the target venue requires anonymity BEFORE submission.
\author{\IEEEauthorblockN{Prajjith Kumar, Thejas K S, Krishnadasa Raja, and S Mahith Shree Vema Reddy}
\IEEEauthorblockA{RV College of Engineering, Bengaluru, India\\
% Insert individually verified institutional email addresses when appropriate.
}}

\begin{document}
\maketitle

\begin{abstract}
A digital signature authenticates a release record, but byte-exact signature
verification does not reveal whether a modified neural-network checkpoint is
still associated with that release. We study a white-box provenance prototype
that combines a SHA-256 dataset commitment, an ML-DSA-65-signed manifest, and a
secret-keyed, Reed--Solomon-coded watermark embedded in convolutional weights.
A trusted auditor compares the extracted 256-bit manifest digest and can retry
extraction after specified channel permutations and positive scalings. Our
central question is how the watermark-training objective trades classifier
accuracy against resistance to an adversarial overwrite. On CIFAR-10 with a
160,234-parameter SmallCNN, five paired training seeds give mean test
accuracies of 84.60\% (unmarked), 82.37\% (binary cross-entropy with RS32,
i.e., 32 parity bytes), and 84.10\% (hinge with RS32). Hinge improves mean
accuracy over the BCE configuration by 1.73 percentage points, but all three
hinge checkpoints tested under a 512-bit same-layer overwrite become rejected
with accuracy losses of 3.38--4.71 points; no successful removal is recorded
for either BCE configuration in the tested suite. Across three watermark
designs, the archive contains 252 checkpoint evaluations per design, including
repeated attacks. All 189 recorded permutation/combined-edit and input-scaling
evaluations verify exactly with retries available, and 778 sampled negative
controls produce no detections. These findings indicate an
accuracy--overwrite tradeoff in the evaluated setting; they are not a proof of
ownership, a general robustness result, or an end-to-end false-accept
guarantee.
\end{abstract}

\begin{IEEEkeywords}
Neural-network watermarking, model provenance, post-quantum signatures,
ML-DSA-65, Reed--Solomon coding, structural obfuscation, adversarial overwrite.
\end{IEEEkeywords}

\section{Introduction}
\label{sec:introduction}
Distributing a trained model raises two related but distinct verification
problems. A signed manifest can authenticate a developer's statement about a
particular release. Once a model is pruned, quantized, fine-tuned, or
reparameterized, however, the checkpoint's bytes may no longer match the
original artifact, and the signature alone cannot say whether the modified
checkpoint still belongs to that release. A mark carried in the weights could
supply additional evidence linking the checkpoint to its signed release without
requiring identical files. Such evidence is useful only if its trust
assumptions and its behavior under model modifications are stated explicitly.

Post-quantum digital signatures are relevant to long-lived release records:
a sufficiently capable quantum computer running Shor's algorithm would threaten
classical schemes based on integer factorization or discrete logarithms
\cite{shor1994algorithms}. NIST standardized the lattice-based ML-DSA signature
scheme in FIPS~204 \cite{nist204}. Replacing a release signature, however, does
not make a neural watermark quantum-secure, and a watermark does not establish
who legally owns a model. The two mechanisms provide different, complementary
checks.

Parameter watermarking is well established \cite{uchida2017embedding}, and
studies of removal attacks and structural obfuscation have shown that a mark
which survives ordinary fine-tuning need not survive more deliberate
modifications \cite{lukas2022sok,yan2023rethinking}. We therefore do not
present a digital signature or a weight projection as new. The open point we
address is narrower and practical. A signed release record and a weight
watermark are different kinds of evidence, and their combination raises two
questions that a deployment must answer together: what the embedding objective
costs the host classifier in clean accuracy, and how the resulting mark fares
under deliberate overwriting. A mark should be judged on both, because a
training objective can improve one at the expense of the other. We measure
them jointly, on paired training seeds and under a common attack suite.

We address three research questions. \emph{(RQ1)} What clean-accuracy cost
accompanies training a full-digest, error-corrected weight mark? \emph{(RQ2)}
How do binary-cross-entropy (BCE) and margin-capped hinge objectives compare
under accuracy-constrained overwriting? \emph{(RQ3)} Which recorded changes
preserve exact recovery, residual evidence, or neither, and how should those
outcomes be interpreted? Our contributions are as follows.
\begin{itemize}
\item We implement a testable release-to-checkpoint association: an ordered
SHA-256 Merkle dataset commitment and ML-DSA-65 manifest signature external to
the model, plus a full 256-bit digest embedded as a keyed, error-corrected
weight projection.
\item We specify a trusted white-box verifier with exact digest comparison,
structural extraction retries, and a distinct residual-agreement verdict. We
state the limited probability model behind the retry-eligibility rule rather
than treating it as a proof of overall false acceptance.
\item We compare BCE/RS32, BCE/RS16, and hinge/RS32 on paired training seeds
and a saved attack suite. The central empirical result is a tradeoff: the
hinge objective gives higher clean accuracy than BCE/RS32, but lower resistance
to the specified same-layer overwrite.
\end{itemize}

\emph{Scope of claims.} We use three terms precisely. \emph{Signature
authentication} means that a manifest verifies under a trusted public key.
\emph{Model association} means that a submitted checkpoint carries a payload
matching that manifest's digest under a pinned extraction profile and a
protected watermark key (Section~\ref{sec:design}). \emph{Proof of ownership}
is not established by either mechanism, and we make no such claim. The rest of
the paper is organized as follows: Section~\ref{sec:related} positions the work,
Section~\ref{sec:design} describes the design and verifier,
Section~\ref{sec:setup} gives the protocol, Section~\ref{sec:results} reports
results, and Section~\ref{sec:discussion} states limitations.

\section{Prior Work and Positioning}
\label{sec:related}
\subsection{Neural Watermarks and Adversarial Evaluation}
Uchida \emph{et al.} introduced weight-regularized neural watermarking, which
uses model parameters and a projection-based extraction rule
\cite{uchida2017embedding}. Our parameter carrier and BCE baseline build on
that line of work; we do not claim to originate the projection mechanism.
Other work uses queries to a model rather than its weights. For example, Adi
\emph{et al.} proposed a black-box watermark based on backdoor behavior
\cite{adi2018backdooring}. Black-box and white-box marks support different
verification settings and cannot be compared using raw detection percentages
alone.

Lukas \emph{et al.} studied the robustness of image-classification
watermarks under removal attacks and emphasized the importance of adversarial
evaluation \cite{lukas2022sok}. Yan \emph{et al.} examined neural structural
obfuscation, including function-preserving transformations and dummy-neuron
constructions \cite{yan2023rethinking}. We evaluate selected compensated
permutations and scalings with a trusted extraction key; our verifier does
\emph{not} address the full structural-obfuscation attack space. We also study
an overwrite that trains a competing mark into the same layer, and we report
the victim's classifier accuracy as a separate outcome rather than folding it
into a detection rate.

\subsection{Post-Quantum Release Authentication}
ML-DSA-65 provides standardized digital signatures for release statements
\cite{nist204}. Post-quantum signing of federated model updates
\cite{li2024federated} and model-bill-of-materials attestation
\cite{campbell2026mbom} address adjacent supply-chain questions. Work on
post-quantum-assisted watermarking of generated media, such as diffusion-model
outputs \cite{hu2026diffusion}, addresses a different carrier and threat
model from classifier parameters.

\subsection{Positioning of This Study}
The cited lines of work address, respectively, weight and black-box watermarks,
the evaluation of their removal resistance, and post-quantum authentication of
statements or supply-chain records. Our contribution lies in their integration
and measurement, not in any one component. We bind a signed manifest to a
separately keyed model-weight association, and we measure how the choice of
embedding objective shifts the balance between clean accuracy and overwrite
resistance. The experiments are not a cryptographic comparison with the
post-quantum systems above, and they are not a head-to-head comparison with
other watermarking schemes.

\section{Design and Verification Model}
\label{sec:design}
\begin{figure*}[!t]
 \centering
 \figslot{figures/system_workflow.png}{1.22in}{Two-lane pipeline: signed release creation (top) and trusted white-box checkpoint verification (bottom).}
 \caption{Proposed release-to-checkpoint association. The ML-DSA-65 signature authenticates the manifest under a trusted public key; the separately keyed watermark is evidence of checkpoint association. Dataset comparison is an optional, distinct check.}
 \label{fig:workflow}
\end{figure*}

\subsection{Actors, Scope, and Authenticated Claims}
A developer prepares a model release and its signed manifest. A trusted auditor
has white-box access to a suspect checkpoint, obtains the developer's public
signing key through an independent channel, and protects a separate secret
watermark key. The auditor pins the extraction profile---carrier layer,
aggregation mode, coding variant, full payload length, and alignment policy---
to trusted release context. The reference implementation can otherwise accept
settings from an unsigned command line or local sidecar; our verification
claims assume the auditor does not allow an adversary to select these settings.
The legacy short-payload override is disabled.

The evaluated remover can modify the weights. Its same-layer overwrite knows
the carrier and trains a competing mark using another key, but is not given
the original watermark key or unrestricted access to an adaptive verification
oracle. We also evaluate unrelated weights, wrong keys, and selected tampered
records. These settings do not cover key compromise, black-box extraction,
distillation, model-width changes, or unlimited adaptive submissions.

Table~\ref{tab:claims} summarizes what each component does and does not
establish. A valid signature establishes authentication of the signed statement
under the trusted public key. A matching model watermark supports association
with that statement under the specified extractor. A Merkle root commits to an
ordered byte representation of a dataset. None of these observations alone
establishes that the committed data were actually used for training, that the
dataset is legally owned, or that a self-declared timestamp proves priority.

\begin{table}[!t]
\centering\caption{Scope of Each Verification Component. None Is a Proof of Ownership.}
\label{tab:claims}\scriptsize\setlength{\tabcolsep}{3pt}
\renewcommand{\arraystretch}{1.1}
\begin{tabular}{@{}>{\raggedright\arraybackslash}p{0.2\columnwidth}>{\raggedright\arraybackslash}p{0.36\columnwidth}>{\raggedright\arraybackslash}p{0.36\columnwidth}@{}}
\toprule
Component & Supports & Does not establish\\
\midrule
ML-DSA-65 signature on the manifest & Authentication of the signed statement under the trusted public key & Any property of a specific checkpoint; model creation or ownership\\
Watermark verdict \textsc{Verified} & Exact 256-bit digest match: association of the weights with the signed manifest under the pinned profile and key & Ownership, training history, or security against key compromise or adaptive adversaries\\
Watermark verdict \textsc{Damaged} & Residual agreement with the expected codeword above a cutoff, without an exact match & Exact recovery; a probability of origin\\
Merkle dataset root & Commitment to an ordered byte representation of a dataset & Actual use in training; legal ownership; priority\\
\bottomrule
\end{tabular}
\end{table}

\subsection{Dataset Commitment and Signed Manifest}
For the ordered CIFAR-10 training split, the implementation hashes each labeled
sample using an index and separate domains for leaf and internal nodes. For
sample index $i$, class label $y_i$, and the raw height--width--channel byte
string $x_i$, the leaf is
\begin{equation}
 h_i=H\!\left(\texttt{0x00}\parallel\operatorname{u64be}(i)
 \parallel\operatorname{u16be}(y_i)\parallel x_i\right),
 \label{eq:leaf}
\end{equation}
where $H$ is SHA-256 and $\parallel$ denotes concatenation. Parent nodes use
$H(\texttt{0x01}\parallel h_L\parallel h_R)$; an unpaired node is promoted
without duplication. This yields a 32-byte root committing to the ordered
50,000-image representation, subject to the assumed collision resistance and
fixed serialization \cite{merkle1987digital,cifar10}.

The root and release metadata are placed in a project-defined deterministic
JSON serialization with sorted keys, compact separators, and ASCII escaping
before UTF-8 encoding. The manifest records the developer identifier, public
key identifier, dataset description, model architecture and parameter count,
payload length, and a UTC timestamp. Let $M$ be the serialized bytes. The
record is signed using ML-DSA-65 \cite{nist204}, and the embedded payload is
\begin{equation}
 p=\operatorname{SHA256}(M)\in\{0,1\}^{256}.
 \label{eq:payload}
\end{equation}
An ML-DSA-65 signature (3,309 bytes) and public key (1,952 bytes) are kept
outside the model weights; the smaller digest is the actual watermark payload.
Signing the manifest and embedding its digest should not be confused with
signing the model weights themselves.

\subsection{Keyed Carrier and Error-Corrected Payload}
We select the final convolutional weight tensor of shape
$128\times64\times3\times3$ (\texttt{features.6.0.weight}).
Averaging its 128 output filters gives a $d=576$-dimensional carrier
$v\in\mathbb{R}^{576}$. A deterministic secret-keyed $\{-1,+1\}$ projection
matrix $X\in\mathbb{R}^{n\times576}$ is obtained through a domain-separated
SHA-256 derivation followed by SHAKE-256 expansion. Its output scores are
$z=Xv$; a score strictly above zero encodes bit one, otherwise bit zero.
Averaging makes this carrier invariant to output-filter permutation in exact
arithmetic, but not inherently invariant to all layer transformations.

The 32-byte payload is systematically encoded using Reed--Solomon coding
over $\mathrm{GF}(256)$ \cite{reed1960polynomial}. RS32 adds 32 parity
bytes, producing $n=512$ embedded bits; RS16 adds 16 parity bytes, producing
$n=384$. Both encode the same 256-bit manifest digest. With $r$ parity
bytes, bounded-distance decoding can correct $e$ erroneous bytes and $f$
known erasures when $2e+f\le r$; confidence-based erasure selection does not
remove that condition.

For codeword bits $b_j\in\{0,1\}$, define $s_j=2b_j-1$. With task
cross-entropy $\mathcal{L}_{\mathrm{cls}}$, we train using
\begin{equation}
 \mathcal{L}=\mathcal{L}_{\mathrm{cls}}+
 \lambda\mathcal{L}_{\mathrm{wm}},
 \label{eq:total}
\end{equation}
where the two investigated watermark penalties are
\begin{align}
 \mathcal{L}_{\mathrm{BCE}}&=\frac1n\sum_{j=1}^n
 \log\!\left(1+e^{-s_jz_j}\right),\label{eq:bce}\\
 \mathcal{L}_{\mathrm{hinge}}&=\frac1n\sum_{j=1}^n
 \max\{0,m-s_jz_j\}.\label{eq:hinge}
\end{align}
At finite signed margins, BCE continues to encourage increasing score
magnitudes. The hinge term ceases contributing a watermark gradient for a bit
once its signed score reaches the prescribed margin $m$, although task loss
and weight decay can still change the weights. This gradient difference
motivates, but does not prove, a connection to the observed overwrite
behavior.

\subsection{Exact Recovery, Structural Retries, and Residual Evidence}
The verifier first checks release authenticity, the public-key identifier,
payload consistency, supported architecture, parameter count, and finite
model tensors. The reported profile assumes restricted PyTorch tensor loading;
a legacy fallback on older environments must not be treated as satisfying
that assumption. For a candidate carrier, the decoder attempts hard RS
decoding, then retries with 2, 4, $\ldots$, $r/2$ low-confidence byte
erasures. Byte confidence is the smallest absolute bit score in the byte.
The first decoder success ends that candidate's erasure schedule, but only a
full 256-bit digest match can yield \textsc{Verified}.

If raw extraction does not match, the enabled verifier may attempt a
key-dependent channel-block assignment, a positive-scale fit, and an
assignment on normalized blocks followed by scaling. The assignment
maximizes a signed-score sum over channel permutations (a linear assignment
problem \cite{kuhn1955hungarian}). The positive-scale
linear program maximizes the worst signed score under a normalization and a
positive lower bound. These are search procedures, not guaranteed inverses
of arbitrary obfuscations. The archived structural outcomes were recorded
with such retries available; they do not provide a matched ablation with the
retry mechanism disabled.

An eligibility test uses Cover's count of strict homogeneous sign patterns
\cite{cover1965geometrical}. For $n$ fixed rows in $C$ dimensions, with a
\emph{uniform independently sampled} target sign vector, the share of
strict patterns realizable by an unrestricted real scale vector is bounded by
\begin{equation}
 B(n,C)=2^{1-n}\sum_{k=0}^{C-1}\binom{n-1}{k}.
 \label{eq:cover}
\end{equation}
With $C=64$, $\log_2 B$ is approximately $-239.71$ for RS32 and $-139.86$
for RS16. These values motivate the implemented retry gate. They do not
bound the full verifier: the actual target is a structured RS codeword,
assignment is target-dependent, and decoding accepts certain incorrect raw
sign patterns. No general false-accept probability for the complete
assignment-plus-decoder pipeline is claimed.

After raw and permitted retries fail exact comparison, the verifier may
report \textsc{Damaged} when the original-order carrier still agrees with
the expected codeword on enough positions. Writing
$\widehat b_j=\mathbf{1}[z_j>0]$, define
\begin{equation}
 k=\sum_{j=1}^n \mathbf{1}[\widehat b_j=b_j].
 \label{eq:agree}
\end{equation}
The cutoffs are $k\ge358$ for RS32 and $k\ge280$ for RS16. Under a
prespecified, one-shot comparison whose hard decisions are independent fair
bits, the corresponding binomial upper tails are about $2^{-64.19}$ and
$2^{-64.21}$. This is a \emph{conditional null calculation}, not an empirical
adaptive-forgery rate or a probability that a model has a particular origin.
Extraction or structural-check failure, or inadequate agreement, leads to
\textsc{Rejected}. Only \textsc{Verified} is an exact digest match.

\section{Experimental Protocol}
\label{sec:setup}
\subsection{Model, Data, and Training}
Experiments use CIFAR-10's 50,000 training and 10,000 test examples
\cite{cifar10} with a width-32 SmallCNN: five convolution--BatchNorm--ReLU
blocks, three max-pooling stages, a dropout-0.3 classifier head, and 160,234
trainable parameters. Training uses random crops with padding four,
horizontal flips, CIFAR-10 normalization, SGD with Nesterov momentum 0.9 and
weight decay $5\times10^{-4}$, batch size 128, an initial learning rate of
0.05, and a 30-epoch cosine schedule. The archived Colab training setup uses
a Tesla T4. The principal watermark configurations set $\lambda=20$, the
final convolution as carrier, and $m=1.5$ for hinge/RS32. We use independent
training seeds 0--4 where available.

Clean accuracy is measured on five baseline, five BCE/RS32, and five
hinge/RS32 checkpoints. BCE/RS16 has three trained checkpoints (seeds 0--2).
The main robustness comparison reuses the first three trained checkpoints
per marked design. It does not constitute 252 independent training trials:
per design, each of three parents produces 84 recorded evaluations across
47 modifications and one clean configuration. Some configurations are
repeated with stochastic attack draws, giving 252 rows per design. The same
release, watermark key, and test set are reused. The implementation and saved
artifacts are archived at an immutable project revision \cite{projectResults}.

\subsection{Attacks, Outcomes, and Utility Constraint}
The attack suite covers Gaussian weight noise, global and carrier-layer
magnitude pruning, FP16 and low-bit quantization, channel permutations,
positive compensating scalings, ordinary and post-pruning fine-tuning, and
same-layer watermark overwrites. The 128- and 512-bit overwrites train a
competing mark on a 20\% CIFAR-10 training subset for 10 epochs, with learning
rate 0.01 and watermark weight $\lambda=20$. Their attacker settings are
shared across the marked designs; this does not match their pre-attack
margins or attacker bit-error rates.

We record \textsc{Verified}, \textsc{Damaged}, and \textsc{Rejected} separately
from test accuracy. For a parent model with clean accuracy $a_0$ and attacked
accuracy $a$ (both expressed as fractions of the test set), an operational
\emph{successful removal} requires
\begin{equation}
 \textsc{Rejected}\quad\text{and}\quad a_0-a\le0.05.
 \label{eq:removal}
\end{equation}
The five-point budget is relative to each parent; it does not impose an
absolute minimum accuracy. Accuracy alone cannot establish preservation of
functionally identical predictions, and a decoded mark can survive even when
the classifier becomes practically unusable. All comparative rates are
descriptive, with parent-model and key reuse explicitly acknowledged.

\section{Results}
\label{sec:results}
The central result is an accuracy--overwrite tradeoff. Hinge/RS32 attains a
higher mean clean accuracy than BCE/RS32 (84.10\% versus 82.37\%) but, under
the 512-bit same-layer overwrite, all three hinge checkpoints are
\textsc{Rejected} within the five-point budget, whereas all six BCE
checkpoints remain \textsc{Damaged}. The subsections below give the paired
accuracy evidence (RQ1), the overwrite evidence (RQ2), and the transformation
and control results (RQ3).

\subsection{RQ1: Clean Accuracy and Carrier Structure}
Table~\ref{tab:clean} reports the paired training results. The five-seed
baseline mean is 84.60\%; BCE/RS32 yields 82.37\%, a mean paired loss of
2.23 percentage points. Hinge/RS32 reaches 84.10\%, a mean paired loss of
0.50 points. Hinge exceeds BCE/RS32 in every one of the five matched seeds,
with paired differences of 1.71, 2.40, 2.17, 0.92, and 1.46 points (mean
1.73). This is a descriptive, five-pair observation, not a population-level
significance claim. The three-seed BCE/RS16 mean of 82.26\% should be compared
with BCE/RS32's matched three-seed mean of 81.98\%, not its five-seed mean.
All 13 main marked training summaries report zero decoded payload errors.

\begin{table}[!t]
\centering\caption{Paired CIFAR-10 Test Accuracy (\%)}
\label{tab:clean}\footnotesize\setlength{\tabcolsep}{3.6pt}
\begin{tabular}{@{}crrrr@{}}
\toprule
Seed & Base & BCE/RS16 & BCE/RS32 & Hinge/RS32\\
\midrule
0 & 86.12 & 82.09 & 83.08 & 84.79\\
1 & 85.32 & 80.75 & 80.33 & 82.73\\
2 & 83.78 & 83.95 & 82.52 & 84.69\\
3 & 81.98 & -- & 82.84 & 83.76\\
4 & 85.78 & -- & 83.07 & 84.53\\
\midrule
Mean & 84.60 & 82.26 & 82.37 & 84.10\\
Paired $\Delta$ (pt) & -- & $-2.81$ & $-2.23$ & $-0.50$\\
\bottomrule
\end{tabular}
\vspace{2pt}
\parbox{0.95\columnwidth}{\scriptsize BCE/RS16 uses seeds 0--2; its paired difference uses those baseline seeds. Other means use seeds 0--4.}
\end{table}

\begin{figure}[!t]
\centering\figslot{figures/paired_accuracy.png}{1.18in}{Paired clean accuracy by seed for baseline, BCE/RS32, hinge/RS32, and the three RS16 seeds.}
\caption{Clean accuracy for the available training seeds. Paired seed comparisons, rather than aggregate means with unequal denominators, are the primary evidence for the loss-function tradeoff.}
\label{fig:acc}
\end{figure}

Averaging the carrier filters reveals an associated change in layer weights.
Across seeds 0--2 the mean-filter norm is approximately 5.05 for BCE/RS32
versus 2.39 for hinge/RS32. The fraction of carrier-layer squared norm in
the shared mean component is approximately 97\% for BCE and 88--89\% for
hinge. Clean minimum absolute projection scores in the saved main runs are
2.24 (BCE/RS32), 2.70 (BCE/RS16), and 1.50 (hinge/RS32). The measurements
suggest a different distribution of embedding effort, but do not by
themselves prove what causes either classification loss or overwrite
susceptibility. A single-seed BCE strength sweep ($\lambda=1,5,20$) yielded
84.67\%, 83.19\%, and 83.08\% clean accuracy, respectively, with minimum
absolute scores 0.02, 1.21, and 2.25; it is exploratory rather than a
multi-seed dose-response study.

\subsection{RQ2: Overwrite Resistance and Classifier Utility}
\begin{table}[!t]
\centering\caption{Recorded Outcomes Over 252 Evaluations per Design}
\label{tab:verdicts}\footnotesize\setlength{\tabcolsep}{3.3pt}
\begin{tabular}{@{}lrrrr@{}}
\toprule
Design & Verified & Damaged & Rejected & Removals\\
\midrule
BCE/RS32 & 231 & 20 & 1 & 0\\
BCE/RS16 & 240 & 12 & 0 & 0\\
Hinge/RS32 & 233 & 13 & 6 & 3\\
\bottomrule
\end{tabular}
\vspace{2pt}
\parbox{0.96\columnwidth}{\scriptsize Includes one clean configuration per seed. Each design uses three independent trained parent checkpoints, not 252 independent models.}
\end{table}

The 252 verdict counts per design are shown in Table~\ref{tab:verdicts}.
The decisive difference appears under the 512-bit same-layer overwrite
(Table~\ref{tab:overwrite}): every tested hinge model is \textsc{Rejected}
with a clean-to-attacked test-accuracy drop between 3.38 and 4.71 points,
within the specified five-point budget. All six BCE models evaluated in the
same overwrite setting (three per coding variant) are instead
\textsc{Damaged}; one BCE/RS32 and two BCE/RS16 trials meet the utility
budget. Thus the observed hinge accuracy benefit coexists with reduced
resistance to this particular watermark-removal setting.

\begin{table}[!t]
\centering\caption{Same-Layer 512-bit Overwrite, Seeds 0--2}
\label{tab:overwrite}\footnotesize\setlength{\tabcolsep}{3pt}
\begin{tabular}{@{}llrrl@{}}
\toprule
Design & Seed & Attacked acc. & Drop (pt) & Verdict\\
\midrule
BCE/RS32 & 0 & 77.79\% & 5.29 & Damaged\\
 & 1 & 74.03\% & 6.30 & Damaged\\
 & 2 & 78.19\% & 4.33 & Damaged\\
BCE/RS16 & 0 & 75.31\% & 6.78 & Damaged\\
 & 1 & 77.35\% & 3.40 & Damaged\\
 & 2 & 80.10\% & 3.85 & Damaged\\
Hinge/RS32 & 0 & 80.08\% & 4.71 & Rejected\\
 & 1 & 79.35\% & 3.38 & Rejected\\
 & 2 & 80.21\% & 4.48 & Rejected\\
\bottomrule
\end{tabular}
\end{table}

\begin{figure*}[!t]
\centering\figslot{figures/robustness_overwrite.png}{1.48in}{Left: stacked counts VERIFIED/DAMAGED/REJECTED for the three designs. Right: per-seed accuracy drop under 512-bit overwrite with verdict marks and five-point budget.}
\caption{Two complementary views of robustness. Left: recorded verdict counts across the full suite, including clean controls. Right: per-parent utility impact of the 512-bit overwrite. Only rejections at or below the five-point drop threshold qualify as successful removals.}
\label{fig:overwrites}
\end{figure*}

The other three hinge rejections arise under two-bit quantization, where
mean test accuracy was only 9.90\%; they are not successful removals under
Eq.~\eqref{eq:removal}. Conversely, a \textsc{Damaged} verdict does not mean
that the original digest was exactly recovered. At the weaker 128-bit
overwrite, BCE/RS16 verified exactly in all three recorded tests.

\subsection{RQ3: Transformations, Error Correction, and Negative Controls}
All 108 recorded evaluations of the specified permutations and combined
permutation/input-scale modification yielded exact matches. Another 81
input-scaling-only evaluations also verified, with unchanged recorded test
accuracy: 189 of 189 in these two families. An additional 81 output-side
rescaling evaluations verified, with absolute measured accuracy changes of
at most 0.07 points. The verifier was run with retries enabled, but the
archive does not identify which candidate succeeded in each case and does
not contain a matched, current-verifier disabled-retry replay. These results
therefore establish recorded exact recovery, not the isolated causal benefit
of each retry.

\begin{figure}[!t]
\centering\figslot{figures/structural_edits.png}{1.12in}{For seven recorded structural configurations, show observed mean clean-relative sign-change rate and the exact verified count (9/9 per design per setting).}
\caption{Selected structural transformations for BCE/RS32, BCE/RS16, and hinge/RS32. Each of seven displayed configurations has nine evaluations per design. Sign-change rate is relative to the clean checkpoint, not agreement with the expected codeword; verification counts are measured with retries enabled.}
\label{fig:structure}
\end{figure}

Exact recovery and useful classifier behavior must be interpreted
separately. For BCE/RS32, 70\% global pruning left mean test accuracy at
16.48\%, and three-bit per-tensor quantization left 24.41\%, even though
all three models verified exactly in each condition. At three bits, ordinary
hard RS decoding corrected 4, 5, and 7 bytes (mean 5.33) with zero erasures;
the average raw sign-change rate relative to clean scores was 1.11\%.
At two bits, BCE/RS32 mean accuracy was 9.53\%; one rejected checkpoint
had 8.58\% accuracy, while the other two were \textsc{Damaged}. These are
examples of recoverability despite severe utility loss, not evidence that
classification was preserved.

The RS32 negative-control study comprises 724 extraction cases, including
600 wrong-key trials, baseline and scrambled checkpoints, unrelated model
weights, and random initialization. Its largest recorded original-order raw
agreement was 305/512, below the 358/512 \textsc{Damaged} cutoff. A separate
set of 54 attacked unmarked-baseline evaluations was also entirely
\textsc{Rejected}. Across these 778 sampled cases, no \textsc{Verified} or
\textsc{Damaged} verdict was observed. Several controls reuse checkpoints,
and the 54 attacked-baseline cases lack the raw-agreement values used in the
724-case histogram. Zero observed detections is not a population-wide false
acceptance guarantee.

\begin{figure}[!t]
\centering\figslot{figures/negative_controls.png}{1.10in}{Histogram of raw agreement for 724 RS32 extraction controls; vertical threshold line at 358/512; observed maximum 305/512.}
\caption{Measured original-order codeword agreement for 724 RS32 extraction controls. The cutoff is 358 of 512; the 54 separately verified attacked baselines are excluded from the distribution.}
\label{fig:controls}
\end{figure}

For context, the archive also contains an exploratory sign-stamping baseline
in which each RS32 code bit controls the sign of one keyed weight after
training. Over 36 equally weighted common attack configurations, this
pipeline yields 26\% exact verification and 61\% exact-or-residual detection,
compared with 86\% and 99\% for the trained projection pipeline. These are
\emph{not} matched-sample estimates of an isolated algorithmic improvement:
the naive version uses two seeds, fewer stochastic draws, and no alignment
retry, whereas the projection version uses three seeds and can retry.

\section{Discussion and Limitations}
\label{sec:discussion}
\subsection{What the Tradeoff Establishes}
The five paired accuracy runs show a consistent clean-accuracy advantage for
the implemented hinge objective relative to the BCE/RS32 setting. The three
same-layer overwrite parents show the converse operational pattern: hinge
models lose both exact recovery and residual detection under the 512-bit
attack at a cost within the preset five-point accuracy budget, whereas the BCE
variants retain residual evidence. The two outcomes therefore pull in opposite
directions in this setting, and a designer should not assume the two objectives
are interchangeable. Reporting exact matching, residual statistical agreement,
and preserved task accuracy separately is more informative than a single
``robustness'' percentage, because each answers a different question.

The weight-carrier measurements are compatible with an explanation in which
BCE keeps increasing shared projection magnitude whereas the hinge penalty
ceases contributing above a target margin. They are not a controlled causal
experiment. The designs differ in trained margins and carrier statistics;
RS16 also changes redundancy and codeword length. A conclusive causal
comparison would match initial margin and watermark strength, attack budgets,
independent keys, and repeated attacker initialization.

\subsection{Security and Reproducibility Boundaries}
The external ML-DSA signature authenticates the manifest under standard
signature assumptions; it neither cryptographically protects the embedding
channel nor proves model creation or ownership. An adversary who learns the
watermark key or can query the verifier adaptively is outside the evaluated
attack model. The signed root covers a reported dataset serialization, not a
certified training history. The timestamp is developer-supplied and lacks an
independent priority witness. A trusted auditor must safely manage the
watermark key and independently authenticate the extraction profile; the
repository's unsigned sidecars are not a substitute for that trust.

Cover's strict fixed-order sign-count bound and the raw-score binomial tail
apply to distinct idealized events. Neither is a bound on the deployed,
target-adaptive assignment plus RS decoder. The 778 clustered negative
controls support a statement about those samples only. No formal
post-quantum reduction is supplied for the weight watermark itself.

Finally, this is one 160,234-parameter architecture on one classification
dataset, with three robustness parent checkpoints per design, one release,
and one watermark key. The current public checkout lacks the original
research watermark key, preventing an independent historical alignment-on/off
replay. The reference implementation and checked-in release signature have
been exercised by the project's validation suite, but a complete third-party
reproduction of the historical marked models would need authorized access
to matching verification material. Larger architectures, multiple releases,
stronger adaptive overwrites, dummy-neuron insertions, distillation,
black-box verification, and longer fine-tuning remain open. Independent
signing-key custody and timestamping would be deployment extensions.

\section{Conclusion}
\label{sec:conclusion}
We examined the interaction of post-quantum-signed release metadata and a
keyed neural-weight watermark as a scoped provenance-association mechanism.
The main empirical result is an accuracy--overwrite tradeoff: hinge/RS32
improves clean accuracy by 1.73 points over BCE/RS32 across five paired seeds,
but in all three tested parent checkpoints it meets the operational removal
criterion under the specified 512-bit same-layer overwrite (\textsc{Rejected}
within the five-point utility budget), which no BCE configuration does in the
tested suite. Exact recovery under selected model transformations and the
absence of detections in sampled controls provide complementary, limited
evidence about the implemented verifier. Future evaluation should test
independent keys and architectures, matched-margin attacks, and stronger
adaptive and structural adversaries. Keeping signature authentication, model
association, and proof of ownership distinct is essential when neural
watermarks are used to interpret signed release records.

% IEEE requires disclosure of AI-generated content. Confirm the precise
% tool/section description and whether the conference has a separate policy.
\section*{Acknowledgment and AI-Assistance Disclosure}
% AUTHOR CHECK: Confirm this acknowledgment with the mentor and coauthors.
% AUTHOR CHECK: Verify the tool/model names below (e.g., "GPT-6") against
% what was actually used; they were carried over unchanged from Version 1.
The authors thank Dr.\ Nivya Muchikel of RV College of Engineering for
mentorship and project guidance. OpenAI ChatGPT
(GPT-6) assisted with initial drafting and organization of this manuscript,
including its abstract, introduction, methods, evaluation, and discussion;
OpenAI Codex assisted with earlier code and evidence audits, and Anthropic
Claude assisted with editing earlier manuscript drafts and with the editorial
revision of this version. The named authors
are responsible for checking all data, citations, and scientific claims before
submission.

\bibliographystyle{IEEEtran}
\bibliography{references}
\end{document}
